\section{Appendix --- Source Code}

The complete source code for Laboratory Work 3.1 is organized below by architectural layer. All source files are available in the project GitHub repository.

\textbf{GitHub Repository:} \url{https://github.com/avremere/es-labs}

% ══════════════════════════════════════════════════════════════════════════
% DRIVER LAYER
% ══════════════════════════════════════════════════════════════════════════

\subsection{Driver Layer: AnalogTempSensor}

\subsubsection{AnalogTempSensor.h}
\begin{lstlisting}[language=C++, caption=AnalogTempSensor.h --- NTC Thermistor Driver Interface]
/**
 * @file AnalogTempSensor.h
 * @brief NTC Thermistor Analog Temperature Sensor Driver Interface
 *
 * Provides a reusable abstraction for reading temperature from an NTC
 * thermistor connected to an analog ADC pin through a voltage divider
 * circuit. Converts raw ADC readings to temperature in degrees Celsius
 * using the Steinhart-Hart Beta parameter equation.
 *
 * Circuit configuration (voltage divider):
 *   VCC --- [R_series] --- ADC_PIN --- [NTC] --- GND
 *
 * The NTC resistance is calculated from the ADC reading and then
 * converted to temperature using:
 *   1/T = 1/T0 + (1/B) * ln(R_ntc / R0)
 *
 * Usage:
 *   AnalogTempSensor ntc(A0, 10000, 10000, 3950);
 *   ntc.init();
 *   float tempC = ntc.readTemperatureC();
 */

#ifndef ANALOG_TEMP_SENSOR_H
#define ANALOG_TEMP_SENSOR_H

#include <Arduino.h>

class AnalogTempSensor {
public:
    AnalogTempSensor(uint8_t adcPin, uint32_t seriesR,
                     uint32_t nominalR, uint16_t betaCoeff,
                     float nominalTempC = 25.0f,
                     uint16_t adcResolution = 10);
    void init();
    uint16_t readRaw();
    float readResistance();
    float readTemperatureC();
    float readTemperatureF();
    bool isValid() const;
    uint16_t getLastRaw() const;
    float getLastTemperatureC() const;

private:
    uint8_t  _adcPin;
    uint32_t _seriesR;
    uint32_t _nominalR;
    uint16_t _betaCoeff;
    float    _nominalTempK;
    uint16_t _adcMax;
    uint16_t _lastRaw;
    float    _lastTempC;
    bool     _valid;
};

#endif // ANALOG_TEMP_SENSOR_H
\end{lstlisting}

\subsubsection{AnalogTempSensor.cpp}
\begin{lstlisting}[language=C++, caption=AnalogTempSensor.cpp --- NTC Thermistor Driver Implementation]
/**
 * @file AnalogTempSensor.cpp
 * @brief NTC Thermistor Analog Temperature Sensor Driver Implementation
 *
 * Implements the ADC-based temperature reading for an NTC thermistor
 * connected through a voltage divider circuit. Uses the Steinhart-Hart
 * Beta parameter equation for resistance-to-temperature conversion.
 */

#include "AnalogTempSensor.h"
#include <math.h>

// Constructor
AnalogTempSensor::AnalogTempSensor(uint8_t adcPin, uint32_t seriesR,
                                   uint32_t nominalR, uint16_t betaCoeff,
                                   float nominalTempC, uint16_t adcResolution)
    : _adcPin(adcPin),
      _seriesR(seriesR),
      _nominalR(nominalR),
      _betaCoeff(betaCoeff),
      _nominalTempK(nominalTempC + 273.15f),
      _adcMax((1 << adcResolution) - 1),
      _lastRaw(0),
      _lastTempC(NAN),
      _valid(false) {}

void AnalogTempSensor::init() {
    pinMode(_adcPin, INPUT);
    _lastRaw  = 0;
    _lastTempC = NAN;
    _valid    = false;
}

uint16_t AnalogTempSensor::readRaw() {
    _lastRaw = analogRead(_adcPin);
    return _lastRaw;
}

float AnalogTempSensor::readResistance() {
    uint16_t adc = readRaw();

    // Guard against division by zero or open/short circuit.
    if (adc == 0 || adc >= _adcMax) {
        _valid = false;
        return -1.0f;
    }

    // Voltage divider: R_ntc = R_series * ADC / (ADC_MAX - ADC)
    float resistance = (float)_seriesR * (float)adc
                       / (float)(_adcMax - adc);
    _valid = true;
    return resistance;
}

float AnalogTempSensor::readTemperatureC() {
    float resistance = readResistance();
    if (resistance < 0.0f) {
        _lastTempC = NAN;
        _valid = false;
        return NAN;
    }

    // Steinhart-Hart Beta equation:
    // 1/T = 1/T0 + (1/B) * ln(R / R0)
    float steinhart = log(resistance / (float)_nominalR);
    steinhart /= (float)_betaCoeff;
    steinhart += 1.0f / _nominalTempK;
    float tempK = 1.0f / steinhart;
    _lastTempC = tempK - 273.15f;
    _valid = true;
    return _lastTempC;
}

float AnalogTempSensor::readTemperatureF() {
    float tempC = readTemperatureC();
    if (isnan(tempC)) return NAN;
    return tempC * 9.0f / 5.0f + 32.0f;
}

bool AnalogTempSensor::isValid() const {
    return _valid;
}

uint16_t AnalogTempSensor::getLastRaw() const {
    return _lastRaw;
}

float AnalogTempSensor::getLastTemperatureC() const {
    return _lastTempC;
}
\end{lstlisting}

\subsection{Driver Layer: DigitalTempSensor}

\subsubsection{DigitalTempSensor.h}
\begin{lstlisting}[language=C++, caption=DigitalTempSensor.h --- DS18B20 OneWire Driver Interface]
/**
 * @file DigitalTempSensor.h
 * @brief DS18B20 Digital Temperature Sensor Driver Interface (OneWire)
 *
 * Provides a reusable abstraction for reading temperature from a Dallas
 * DS18B20 digital temperature sensor using the OneWire protocol.
 *
 * Circuit configuration:
 *   VCC (5V) --+-- DS18B20 VDD
 *              |
 *           [4.7K]  (pull-up resistor)
 *              |
 *   DATA_PIN --+-- DS18B20 DQ
 *   GND -------+-- DS18B20 GND
 */

#ifndef DIGITAL_TEMP_SENSOR_H
#define DIGITAL_TEMP_SENSOR_H

#include <Arduino.h>
#include <OneWire.h>
#include <DallasTemperature.h>

class DigitalTempSensor {
public:
    DigitalTempSensor(uint8_t dataPin,
                      uint8_t resolution = 10);
    bool init();
    void requestConversion();
    bool isConversionComplete();
    float readTemperatureC();
    float getLastTemperatureC() const;
    float readTemperatureF();
    bool isConnected() const;
    bool isValid() const;
    uint8_t getDeviceCount() const;

private:
    OneWire           _oneWire;
    DallasTemperature _sensors;
    uint8_t           _resolution;
    float             _lastTempC;
    bool              _connected;
    bool              _valid;
    uint8_t           _deviceCount;
};

#endif // DIGITAL_TEMP_SENSOR_H
\end{lstlisting}

\subsubsection{DigitalTempSensor.cpp}
\begin{lstlisting}[language=C++, caption=DigitalTempSensor.cpp --- DS18B20 OneWire Driver Implementation]
/**
 * @file DigitalTempSensor.cpp
 * @brief DS18B20 Digital Temperature Sensor Driver Implementation
 *
 * Implements OneWire-based temperature reading from the DS18B20 sensor.
 * Uses the DallasTemperature library for protocol handling and provides
 * both blocking and non-blocking read modes.
 */

#include "DigitalTempSensor.h"

static const float DISCONNECTED_TEMP = -127.0f;
static const float POWER_ON_RESET_TEMP = 85.0f;

DigitalTempSensor::DigitalTempSensor(uint8_t dataPin,
                                     uint8_t resolution)
    : _oneWire(dataPin),
      _sensors(&_oneWire),
      _resolution(resolution),
      _lastTempC(NAN),
      _connected(false),
      _valid(false),
      _deviceCount(0) {}

bool DigitalTempSensor::init() {
    _sensors.begin();
    _deviceCount = _sensors.getDeviceCount();
    if (_deviceCount == 0) {
        _connected = false;
        return false;
    }
    _connected = true;
    _sensors.setResolution(_resolution);
    _sensors.setWaitForConversion(false);
    return true;
}

void DigitalTempSensor::requestConversion() {
    if (_connected) {
        _sensors.requestTemperatures();
    }
}

bool DigitalTempSensor::isConversionComplete() {
    if (!_connected) return false;
    return _sensors.isConversionComplete();
}

float DigitalTempSensor::readTemperatureC() {
    if (!_connected) {
        _valid = false;
        _lastTempC = NAN;
        return NAN;
    }
    _sensors.setWaitForConversion(true);
    _sensors.requestTemperatures();
    _sensors.setWaitForConversion(false);

    float tempC = _sensors.getTempCByIndex(0);

    if (tempC == DEVICE_DISCONNECTED_C ||
        tempC <= DISCONNECTED_TEMP) {
        _valid = false;
        _lastTempC = NAN;
        return NAN;
    }
    if (tempC == POWER_ON_RESET_TEMP) {
        _valid = false;
        return _lastTempC;
    }
    _valid = true;
    _lastTempC = tempC;
    return _lastTempC;
}

float DigitalTempSensor::getLastTemperatureC() const {
    return _lastTempC;
}

float DigitalTempSensor::readTemperatureF() {
    float tempC = readTemperatureC();
    if (isnan(tempC)) return NAN;
    return tempC * 9.0f / 5.0f + 32.0f;
}

bool DigitalTempSensor::isConnected() const {
    return _connected;
}

bool DigitalTempSensor::isValid() const {
    return _valid;
}

uint8_t DigitalTempSensor::getDeviceCount() const {
    return _deviceCount;
}
\end{lstlisting}

% ══════════════════════════════════════════════════════════════════════════
% SERVICE LAYER
% ══════════════════════════════════════════════════════════════════════════

\subsection{Service Layer: ThresholdAlert}

\subsubsection{ThresholdAlert.h}
\begin{lstlisting}[language=C++, caption=ThresholdAlert.h --- Hysteresis-Based Threshold Alert Interface]
/**
 * @file ThresholdAlert.h
 * @brief Hysteresis-Based Threshold Alert Module Interface
 *
 * 4-state FSM: NORMAL -> DEBOUNCE_HIGH -> ALERT -> DEBOUNCE_LOW -> NORMAL
 * Each transition requires N consecutive confirmations (debounce).
 * Hysteresis band: [lowThreshold, highThreshold]
 */

#ifndef THRESHOLD_ALERT_H
#define THRESHOLD_ALERT_H

#include <Arduino.h>

enum AlertState {
    ALERT_NORMAL,
    ALERT_DEBOUNCE_HIGH,
    ALERT_ACTIVE,
    ALERT_DEBOUNCE_LOW
};

class ThresholdAlert {
public:
    ThresholdAlert(float highThreshold, float lowThreshold,
                   uint8_t debounceCount = 5);
    void init();
    AlertState update(float value);
    AlertState getState() const;
    bool isAlertActive() const;
    bool isDebouncing() const;
    uint8_t getDebounceCounter() const;
    const char* getStateString() const;
    void setThresholds(float high, float low);
    void setDebounceCount(uint8_t count);
    float getHighThreshold() const;
    float getLowThreshold() const;

private:
    float      _highThreshold;
    float      _lowThreshold;
    uint8_t    _debounceMax;
    uint8_t    _debounceCounter;
    AlertState _state;
};

#endif // THRESHOLD_ALERT_H
\end{lstlisting}

\subsubsection{ThresholdAlert.cpp}
\begin{lstlisting}[language=C++, caption=ThresholdAlert.cpp --- Hysteresis-Based Threshold Alert Implementation]
/**
 * @file ThresholdAlert.cpp
 * @brief Hysteresis-Based Threshold Alert Module Implementation
 *
 * Implements the 4-state FSM for threshold detection with hysteresis
 * and debounce filtering.
 */

#include "ThresholdAlert.h"

static const char* STATE_NAMES[] = {
    "NORMAL", "DEBOUNCE_HIGH", "ALERT", "DEBOUNCE_LOW"
};

ThresholdAlert::ThresholdAlert(float highThreshold,
                               float lowThreshold,
                               uint8_t debounceCount)
    : _highThreshold(highThreshold),
      _lowThreshold(lowThreshold),
      _debounceMax(debounceCount),
      _debounceCounter(0),
      _state(ALERT_NORMAL) {}

void ThresholdAlert::init() {
    _state = ALERT_NORMAL;
    _debounceCounter = 0;
}

AlertState ThresholdAlert::update(float value) {
    switch (_state) {
    case ALERT_NORMAL:
        if (value > _highThreshold) {
            _state = ALERT_DEBOUNCE_HIGH;
            _debounceCounter = 1;
        }
        break;
    case ALERT_DEBOUNCE_HIGH:
        if (value > _highThreshold) {
            _debounceCounter++;
            if (_debounceCounter >= _debounceMax) {
                _state = ALERT_ACTIVE;
                _debounceCounter = 0;
            }
        } else {
            _state = ALERT_NORMAL;
            _debounceCounter = 0;
        }
        break;
    case ALERT_ACTIVE:
        if (value < _lowThreshold) {
            _state = ALERT_DEBOUNCE_LOW;
            _debounceCounter = 1;
        }
        break;
    case ALERT_DEBOUNCE_LOW:
        if (value < _lowThreshold) {
            _debounceCounter++;
            if (_debounceCounter >= _debounceMax) {
                _state = ALERT_NORMAL;
                _debounceCounter = 0;
            }
        } else {
            _state = ALERT_ACTIVE;
            _debounceCounter = 0;
        }
        break;
    }
    return _state;
}

AlertState ThresholdAlert::getState() const { return _state; }

bool ThresholdAlert::isAlertActive() const {
    return _state == ALERT_ACTIVE;
}

bool ThresholdAlert::isDebouncing() const {
    return (_state == ALERT_DEBOUNCE_HIGH) ||
           (_state == ALERT_DEBOUNCE_LOW);
}

uint8_t ThresholdAlert::getDebounceCounter() const {
    return _debounceCounter;
}

const char* ThresholdAlert::getStateString() const {
    return STATE_NAMES[_state];
}

void ThresholdAlert::setThresholds(float highThreshold,
                                   float lowThreshold) {
    _highThreshold = highThreshold;
    _lowThreshold  = lowThreshold;
    _state = ALERT_NORMAL;
    _debounceCounter = 0;
}

void ThresholdAlert::setDebounceCount(uint8_t count) {
    _debounceMax = count;
    _debounceCounter = 0;
}

float ThresholdAlert::getHighThreshold() const {
    return _highThreshold;
}

float ThresholdAlert::getLowThreshold() const {
    return _lowThreshold;
}
\end{lstlisting}

% ══════════════════════════════════════════════════════════════════════════
% APPLICATION LAYER: SHARED DATA
% ══════════════════════════════════════════════════════════════════════════

\subsection{Application Layer: Shared Data and Synchronization}

\subsubsection{sensor\_data.h}
\begin{lstlisting}[language=C++, caption=sensor\_data.h --- Shared Data Structures and Pin Mapping]
/**
 * @file sensor_data.h
 * @brief Lab 3.1 -- Shared Data Structures and Synchronization Primitives
 *
 * Declares the inter-task shared data structures, FreeRTOS synchronization
 * handles, hardware pin mapping, and application constants.
 */

#ifndef SENSOR_DATA_H
#define SENSOR_DATA_H

#include <Arduino.h>
#include <Arduino_FreeRTOS.h>
#include <semphr.h>
#include "ThresholdAlert.h"

// === Hardware Pin Mapping (Arduino Mega 2560) ===
static const uint8_t PIN_ANALOG_SENSOR  = A0;
static const uint8_t PIN_DIGITAL_SENSOR = 2;
static const uint8_t PIN_LED_GREEN      = 8;
static const uint8_t PIN_LED_RED        = 9;
static const uint8_t PIN_LED_YELLOW     = 10;
static const uint8_t LCD_I2C_ADDRESS    = 0x27;
static const uint8_t LCD_COLS           = 16;
static const uint8_t LCD_ROWS           = 2;

// === NTC Thermistor Parameters ===
static const uint32_t NTC_SERIES_RESISTANCE  = 10000;
static const uint32_t NTC_NOMINAL_RESISTANCE = 10000;
static const uint16_t NTC_BETA_COEFFICIENT   = 3950;
static const float    NTC_NOMINAL_TEMP_C     = 25.0f;

// === DS18B20 Parameters ===
static const uint8_t DS18B20_RESOLUTION = 10;

// === Threshold Alert Parameters ===
static const float   ANALOG_THRESHOLD_HIGH  = 30.0f;
static const float   ANALOG_THRESHOLD_LOW   = 28.0f;
static const float   DIGITAL_THRESHOLD_HIGH = 30.0f;
static const float   DIGITAL_THRESHOLD_LOW  = 28.0f;
static const uint8_t ALERT_DEBOUNCE_COUNT   = 5;

// === FreeRTOS Task Configuration ===
static const uint32_t TASK_ACQUISITION_PERIOD_MS  = 50;
static const UBaseType_t TASK_ACQUISITION_PRIORITY = 3;
static const configSTACK_DEPTH_TYPE TASK_ACQUISITION_STACK = 256;

static const uint32_t TASK_CONDITIONING_PERIOD_MS  = 100;
static const UBaseType_t TASK_CONDITIONING_PRIORITY = 2;
static const configSTACK_DEPTH_TYPE TASK_CONDITIONING_STACK = 256;

static const uint32_t TASK_DISPLAY_PERIOD_MS  = 500;
static const UBaseType_t TASK_DISPLAY_PRIORITY = 1;
static const configSTACK_DEPTH_TYPE TASK_DISPLAY_STACK = 512;

// === Shared Data Structures ===
typedef struct {
    uint16_t   analogRaw;
    float      analogResistance;
    float      analogTempC;
    bool       analogValid;
    float      digitalTempC;
    bool       digitalValid;
    uint32_t   readingCount;
    TickType_t timestamp;
} SensorReadings_t;

typedef struct {
    AlertState analogAlertState;
    uint8_t    analogDebounceCount;
    float      analogFilteredTemp;
    AlertState digitalAlertState;
    uint8_t    digitalDebounceCount;
    float      digitalFilteredTemp;
    uint32_t   analogAlertCount;
    uint32_t   digitalAlertCount;
    uint32_t   conditioningCycles;
} AlertStatus_t;

// === External Declarations ===
extern SensorReadings_t g_sensorData;
extern AlertStatus_t    g_alertData;
extern SemaphoreHandle_t xNewReadingSemaphore;
extern SemaphoreHandle_t xSensorMutex;

void sensorDataInit();

#endif // SENSOR_DATA_H
\end{lstlisting}

\subsubsection{sensor\_data.cpp}
\begin{lstlisting}[language=C++, caption=sensor\_data.cpp --- Shared State Definitions]
/**
 * @file sensor_data.cpp
 * @brief Lab 3.1 -- Shared State Definitions and Synchronization Init
 */

#include "sensor_data.h"

SensorReadings_t g_sensorData = {
    0, 0.0f, 0.0f, false, 0.0f, false, 0, 0
};

AlertStatus_t g_alertData = {
    ALERT_NORMAL, 0, 0.0f,
    ALERT_NORMAL, 0, 0.0f,
    0, 0, 0
};

SemaphoreHandle_t xNewReadingSemaphore = NULL;
SemaphoreHandle_t xSensorMutex         = NULL;

void sensorDataInit() {
    xNewReadingSemaphore = xSemaphoreCreateBinary();
    xSensorMutex = xSemaphoreCreateMutex();
}
\end{lstlisting}

% ══════════════════════════════════════════════════════════════════════════
% APPLICATION LAYER: TASKS
% ══════════════════════════════════════════════════════════════════════════

\subsection{Application Layer: Task 1 --- Sensor Acquisition}

\subsubsection{task\_acquisition.h}
\begin{lstlisting}[language=C++, caption=task\_acquisition.h --- Sensor Acquisition Task Interface]
/**
 * @file task_acquisition.h
 * @brief Lab 3.1 -- Sensor Acquisition Task Interface (Task 1)
 *
 * Periodically reads both the analog NTC thermistor (via ADC) and
 * the digital DS18B20 sensor (via OneWire) at 50 ms intervals.
 */

#ifndef TASK_ACQUISITION_H
#define TASK_ACQUISITION_H

#include <Arduino_FreeRTOS.h>

void vTaskAcquisition(void *pvParameters);

#endif // TASK_ACQUISITION_H
\end{lstlisting}

\subsubsection{task\_acquisition.cpp}
\begin{lstlisting}[language=C++, caption=task\_acquisition.cpp --- Sensor Acquisition Task Implementation]
/**
 * @file task_acquisition.cpp
 * @brief Lab 3.1 -- Sensor Acquisition Task Implementation (Task 1)
 *
 * Reads NTC thermistor via ADC and DS18B20 via OneWire at 50 ms
 * intervals using vTaskDelayUntil() for drift-free timing.
 */

#include "task_acquisition.h"
#include "sensor_data.h"
#include "AnalogTempSensor.h"
#include "DigitalTempSensor.h"

static AnalogTempSensor s_ntcSensor(
    PIN_ANALOG_SENSOR, NTC_SERIES_RESISTANCE,
    NTC_NOMINAL_RESISTANCE, NTC_BETA_COEFFICIENT,
    NTC_NOMINAL_TEMP_C);

static DigitalTempSensor s_ds18b20(
    PIN_DIGITAL_SENSOR, DS18B20_RESOLUTION);

void vTaskAcquisition(void *pvParameters) {
    (void)pvParameters;

    s_ntcSensor.init();
    bool ds18b20Found = s_ds18b20.init();

    if (ds18b20Found) {
        s_ds18b20.requestConversion();
    }

    bool firstConversion = true;
    TickType_t xLastWakeTime = xTaskGetTickCount();
    const TickType_t xPeriod =
        pdMS_TO_TICKS(TASK_ACQUISITION_PERIOD_MS);

    uint16_t rawAdc;
    float resistance, analogTemp, digitalTemp;
    bool analogOk, digitalOk;

    for (;;) {
        vTaskDelayUntil(&xLastWakeTime, xPeriod);

        // Read analog sensor (NTC thermistor)
        analogTemp = s_ntcSensor.readTemperatureC();
        rawAdc     = s_ntcSensor.getLastRaw();
        resistance = s_ntcSensor.readResistance();
        analogOk   = s_ntcSensor.isValid();

        // Read digital sensor (DS18B20)
        if (ds18b20Found) {
            if (firstConversion) {
                while (!s_ds18b20.isConversionComplete()) {
                    vTaskDelay(pdMS_TO_TICKS(10));
                }
                firstConversion = false;
            }
            if (s_ds18b20.isConversionComplete()) {
                digitalTemp = s_ds18b20.readTemperatureC();
                digitalOk   = s_ds18b20.isValid();
            } else {
                digitalTemp = s_ds18b20.getLastTemperatureC();
                digitalOk   = s_ds18b20.isValid();
            }
            s_ds18b20.requestConversion();
        } else {
            digitalTemp = NAN;
            digitalOk   = false;
        }

        // Write to shared data under mutex
        if (xSemaphoreTake(xSensorMutex,
                           pdMS_TO_TICKS(10)) == pdTRUE) {
            g_sensorData.analogRaw        = rawAdc;
            g_sensorData.analogResistance = resistance;
            g_sensorData.analogTempC      = analogTemp;
            g_sensorData.analogValid      = analogOk;
            g_sensorData.digitalTempC     = digitalTemp;
            g_sensorData.digitalValid     = digitalOk;
            g_sensorData.readingCount++;
            g_sensorData.timestamp = xTaskGetTickCount();
            xSemaphoreGive(xSensorMutex);
        }

        // Notify Task 2
        xSemaphoreGive(xNewReadingSemaphore);
    }
}
\end{lstlisting}

\subsection{Application Layer: Task 2 --- Signal Conditioning}

\subsubsection{task\_conditioning.h}
\begin{lstlisting}[language=C++, caption=task\_conditioning.h --- Signal Conditioning Task Interface]
/**
 * @file task_conditioning.h
 * @brief Lab 3.1 -- Signal Conditioning & Alert Task Interface (Task 2)
 *
 * Event-driven task: wakes on semaphore from Task 1, applies
 * hysteresis-based threshold detection and controls alert LEDs.
 */

#ifndef TASK_CONDITIONING_H
#define TASK_CONDITIONING_H

#include <Arduino_FreeRTOS.h>

void vTaskConditioning(void *pvParameters);

#endif // TASK_CONDITIONING_H
\end{lstlisting}

\subsubsection{task\_conditioning.cpp}
\begin{lstlisting}[language=C++, caption=task\_conditioning.cpp --- Signal Conditioning Task Implementation]
/**
 * @file task_conditioning.cpp
 * @brief Lab 3.1 -- Signal Conditioning & Alert Task (Task 2)
 *
 * Applies threshold detection with hysteresis and debounce to both
 * analog and digital temperature readings. Controls LED indicators.
 */

#include "task_conditioning.h"
#include "sensor_data.h"
#include "ThresholdAlert.h"
#include "Led.h"

static ThresholdAlert s_analogAlert(
    ANALOG_THRESHOLD_HIGH, ANALOG_THRESHOLD_LOW,
    ALERT_DEBOUNCE_COUNT);

static ThresholdAlert s_digitalAlert(
    DIGITAL_THRESHOLD_HIGH, DIGITAL_THRESHOLD_LOW,
    ALERT_DEBOUNCE_COUNT);

static Led s_greenLed(PIN_LED_GREEN);
static Led s_redLed(PIN_LED_RED);
static Led s_yellowLed(PIN_LED_YELLOW);

void vTaskConditioning(void *pvParameters) {
    (void)pvParameters;

    s_analogAlert.init();
    s_digitalAlert.init();
    s_greenLed.init();
    s_redLed.init();
    s_yellowLed.init();

    s_greenLed.turnOn();
    s_redLed.turnOff();
    s_yellowLed.turnOff();

    float analogTemp, digitalTemp;
    bool analogValid, digitalValid;
    AlertState prevAnalogState  = ALERT_NORMAL;
    AlertState prevDigitalState = ALERT_NORMAL;

    for (;;) {
        // Wait for new sensor data
        if (xSemaphoreTake(xNewReadingSemaphore,
                           portMAX_DELAY) != pdTRUE) {
            continue;
        }

        // Read sensor data under mutex
        if (xSemaphoreTake(xSensorMutex,
                           pdMS_TO_TICKS(10)) == pdTRUE) {
            analogTemp   = g_sensorData.analogTempC;
            digitalTemp  = g_sensorData.digitalTempC;
            analogValid  = g_sensorData.analogValid;
            digitalValid = g_sensorData.digitalValid;
            xSemaphoreGive(xSensorMutex);
        } else {
            continue;
        }

        // Apply threshold detection
        AlertState analogState  = ALERT_NORMAL;
        AlertState digitalState = ALERT_NORMAL;

        if (analogValid && !isnan(analogTemp)) {
            analogState = s_analogAlert.update(analogTemp);
        }
        if (digitalValid && !isnan(digitalTemp)) {
            digitalState = s_digitalAlert.update(digitalTemp);
        }

        // Write alert status under mutex
        if (xSemaphoreTake(xSensorMutex,
                           pdMS_TO_TICKS(10)) == pdTRUE) {
            g_alertData.analogAlertState    = analogState;
            g_alertData.analogDebounceCount =
                s_analogAlert.getDebounceCounter();
            g_alertData.analogFilteredTemp  = analogTemp;
            g_alertData.digitalAlertState    = digitalState;
            g_alertData.digitalDebounceCount =
                s_digitalAlert.getDebounceCounter();
            g_alertData.digitalFilteredTemp  = digitalTemp;

            if (analogState == ALERT_ACTIVE &&
                prevAnalogState != ALERT_ACTIVE) {
                g_alertData.analogAlertCount++;
            }
            if (digitalState == ALERT_ACTIVE &&
                prevDigitalState != ALERT_ACTIVE) {
                g_alertData.digitalAlertCount++;
            }
            g_alertData.conditioningCycles++;
            xSemaphoreGive(xSensorMutex);
        }

        prevAnalogState  = analogState;
        prevDigitalState = digitalState;

        // Update LED indicators
        bool anyAlert = (analogState == ALERT_ACTIVE) ||
                        (digitalState == ALERT_ACTIVE);

        if (anyAlert) s_greenLed.turnOff();
        else          s_greenLed.turnOn();

        if (analogState == ALERT_ACTIVE) {
            s_redLed.turnOn();
        } else if (analogState == ALERT_DEBOUNCE_HIGH ||
                   analogState == ALERT_DEBOUNCE_LOW) {
            s_redLed.toggle();  // Blink during debounce
        } else {
            s_redLed.turnOff();
        }

        if (digitalState == ALERT_ACTIVE) {
            s_yellowLed.turnOn();
        } else if (digitalState == ALERT_DEBOUNCE_HIGH ||
                   digitalState == ALERT_DEBOUNCE_LOW) {
            s_yellowLed.toggle();
        } else {
            s_yellowLed.turnOff();
        }
    }
}
\end{lstlisting}

\subsection{Application Layer: Task 3 --- Display and Reporting}

\subsubsection{task\_display.h}
\begin{lstlisting}[language=C++, caption=task\_display.h --- Display and Reporting Task Interface]
/**
 * @file task_display.h
 * @brief Lab 3.1 -- Display & Reporting Task Interface (Task 3)
 *
 * Periodically updates the LCD display and prints structured
 * STDIO reports to the serial terminal at 500 ms intervals.
 */

#ifndef TASK_DISPLAY_H
#define TASK_DISPLAY_H

#include <Arduino_FreeRTOS.h>

void vTaskDisplay(void *pvParameters);

#endif // TASK_DISPLAY_H
\end{lstlisting}

\subsubsection{task\_display.cpp}
\begin{lstlisting}[language=C++, caption=task\_display.cpp --- Display and Reporting Task Implementation]
/**
 * @file task_display.cpp
 * @brief Lab 3.1 -- Display & Reporting Task Implementation (Task 3)
 *
 * Every 500 ms, reads sensor data and alert status, updates the LCD
 * with a summary, and prints a structured STDIO report every 2 seconds.
 */

#include "task_display.h"
#include "sensor_data.h"
#include "LcdDisplay.h"
#include <stdio.h>
#include <string.h>
#include <math.h>

static LcdDisplay s_lcd(LCD_I2C_ADDRESS, LCD_COLS, LCD_ROWS);

static void formatTemp(char *buf, size_t len, float temp) {
    if (isnan(temp)) strncpy(buf, " -- ", len);
    else             dtostrf(temp, 4, 1, buf);
}

static const char* alertFullLabel(AlertState state) {
    switch (state) {
        case ALERT_NORMAL:        return "NORMAL";
        case ALERT_DEBOUNCE_HIGH: return "DEBOUNCE_HIGH";
        case ALERT_ACTIVE:        return "ALERT";
        case ALERT_DEBOUNCE_LOW:  return "DEBOUNCE_LOW";
        default:                  return "UNKNOWN";
    }
}

void vTaskDisplay(void *pvParameters) {
    (void)pvParameters;

    s_lcd.init();
    s_lcd.backlight(true);
    s_lcd.showTwoLines("Lab 3.1 Sensors", "Initializing...");
    vTaskDelay(pdMS_TO_TICKS(1000));

    TickType_t xLastWakeTime = xTaskGetTickCount();
    const TickType_t xPeriod =
        pdMS_TO_TICKS(TASK_DISPLAY_PERIOD_MS);

    uint8_t displayCycle = 0;
    uint32_t reportNumber = 0;
    static const uint8_t REPORT_INTERVAL = 4;

    SensorReadings_t localSensor;
    AlertStatus_t    localAlert;
    char line0[17], line1[17];

    for (;;) {
        vTaskDelayUntil(&xLastWakeTime, xPeriod);
        displayCycle++;

        // Read shared data under mutex
        if (xSemaphoreTake(xSensorMutex,
                           pdMS_TO_TICKS(20)) == pdTRUE) {
            memcpy(&localSensor, &g_sensorData,
                   sizeof(SensorReadings_t));
            memcpy(&localAlert, &g_alertData,
                   sizeof(AlertStatus_t));
            xSemaphoreGive(xSensorMutex);
        } else {
            continue;
        }

        // Update LCD display
        uint8_t page = (displayCycle / 4) % 2;
        if (page == 0) {
            char aT[8], dT[8];
            formatTemp(aT, sizeof(aT),
                       localSensor.analogTempC);
            formatTemp(dT, sizeof(dT),
                       localSensor.digitalTempC);
            snprintf(line0, sizeof(line0),
                     "A:%sC D:%sC", aT, dT);

            bool aA = (localAlert.analogAlertState
                       == ALERT_ACTIVE);
            bool dA = (localAlert.digitalAlertState
                       == ALERT_ACTIVE);
            if (aA && dA)
                snprintf(line1, sizeof(line1), "ALERT: A+D");
            else if (aA)
                snprintf(line1, sizeof(line1), "ALERT: Analog");
            else if (dA)
                snprintf(line1, sizeof(line1), "ALERT: Digital");
            else
                snprintf(line1, sizeof(line1), "System: OK");
        } else {
            snprintf(line0, sizeof(line0), "AH:%.0f AL:%.0f C",
                     ANALOG_THRESHOLD_HIGH,
                     ANALOG_THRESHOLD_LOW);
            snprintf(line1, sizeof(line1), "DH:%.0f DL:%.0f C",
                     DIGITAL_THRESHOLD_HIGH,
                     DIGITAL_THRESHOLD_LOW);
        }
        s_lcd.showTwoLines(line0, line1);

        // Structured STDIO report (every 2 seconds)
        if ((displayCycle % REPORT_INTERVAL) == 0) {
            reportNumber++;
            char aTStr[8], dTStr[8];
            formatTemp(aTStr, sizeof(aTStr),
                       localSensor.analogTempC);
            formatTemp(dTStr, sizeof(dTStr),
                       localSensor.digitalTempC);

            printf("\r\n");
            printf("=== SENSOR REPORT #%lu ===\r\n",
                   (unsigned long)reportNumber);
            printf("Analog:  %s C (raw=%u, R=%.0f) [%s]\r\n",
                   aTStr, localSensor.analogRaw,
                   localSensor.analogResistance,
                   alertFullLabel(
                       localAlert.analogAlertState));
            printf("Digital: %s C [%s]\r\n",
                   dTStr,
                   alertFullLabel(
                       localAlert.digitalAlertState));
            printf("Alerts: A=%lu D=%lu Cycles=%lu\r\n",
                   (unsigned long)
                       localAlert.analogAlertCount,
                   (unsigned long)
                       localAlert.digitalAlertCount,
                   (unsigned long)
                       localAlert.conditioningCycles);
            printf("========================\r\n");
        }
    }
}
\end{lstlisting}

% ══════════════════════════════════════════════════════════════════════════
% APPLICATION LAYER: ENTRY POINT
% ══════════════════════════════════════════════════════════════════════════

\subsection{Application Layer: Lab 3.1 Entry Point}

\subsubsection{lab3\_1\_main.h}
\begin{lstlisting}[language=C++, caption=lab3\_1\_main.h --- Lab 3.1 Entry Point Interface]
/**
 * @file lab3_1_main.h
 * @brief Lab 3.1 Entry Point Interface
 *
 * Declares setup and loop for the dual-sensor temperature
 * monitoring system with FreeRTOS preemptive scheduling.
 */

#ifndef LAB3_1_MAIN_H
#define LAB3_1_MAIN_H

void lab3_1Setup();
void lab3_1Loop();

#endif // LAB3_1_MAIN_H
\end{lstlisting}

\subsubsection{lab3\_1\_main.cpp}
\begin{lstlisting}[language=C++, caption=lab3\_1\_main.cpp --- Lab 3.1 Entry Point Implementation]
/**
 * @file lab3_1_main.cpp
 * @brief Lab 3.1 -- Dual-Sensor Temperature Monitoring
 *
 * Initializes STDIO, creates synchronization primitives, spawns
 * three FreeRTOS tasks, and starts the scheduler.
 */

#include "lab3_1_main.h"
#include "sensor_data.h"
#include "task_acquisition.h"
#include "task_conditioning.h"
#include "task_display.h"

#include <Arduino.h>
#include <Arduino_FreeRTOS.h>
#include <stdio.h>
#include "StdioSerial.h"

void lab3_1Setup() {
    stdioSerialInit(9600);

    printf("\r\n");
    printf("========================================\r\n");
    printf("  Lab 3.1 -- Dual-Sensor Temp Monitor\r\n");
    printf("  FreeRTOS | Tasks: 3 | Tick: %u Hz\r\n",
           (unsigned int)configTICK_RATE_HZ);
    printf("========================================\r\n");
    printf("Sensors: NTC(A0) + DS18B20(D2)\r\n");
    printf("Thresholds: HIGH=%.1f LOW=%.1f\r\n",
           ANALOG_THRESHOLD_HIGH, ANALOG_THRESHOLD_LOW);
    printf("Debounce: %u confirmations\r\n",
           (unsigned int)ALERT_DEBOUNCE_COUNT);
    printf("========================================\r\n\r\n");

    sensorDataInit();

    xTaskCreate(vTaskAcquisition, "Acquire",
        TASK_ACQUISITION_STACK, NULL,
        TASK_ACQUISITION_PRIORITY, NULL);

    xTaskCreate(vTaskConditioning, "Cond",
        TASK_CONDITIONING_STACK, NULL,
        TASK_CONDITIONING_PRIORITY, NULL);

    xTaskCreate(vTaskDisplay, "Display",
        TASK_DISPLAY_STACK, NULL,
        TASK_DISPLAY_PRIORITY, NULL);
}

void lab3_1Loop() {
    // All logic runs in FreeRTOS tasks.
}
\end{lstlisting}
